% =============================================================================
% Appendix G  The Descriptor Construction in Full
% Cited from sec:descriptor-layouts to sec:descriptor-cost.
% Plan: ThesisPlanning.md §6 (G); headline contribution C.
%
% CONTENT  The block rows (PDE and coupled-field chains included); the padded
%          block count; the proof of thm:kernel-equivalence; the term IR and its
%          compilation; the proof of prop:descriptor-cost; block rescaling.
% SOURCES  Planning.md §3.4, §3.7 ("Descriptor"), §6.4, App. A-2;
%          pihm/src/pihm/assemble.py (descriptor mode), ir.py, circuits/lcu.py.
% =============================================================================

\subsection{Block rows in full}
\label{app:descriptor-rows}
% LAND: every block row of A_desc, including the source row
%   (- M_r D0 f_hat^(0)), the padding rows (c_pad f_hat^(b) = 0; the padding
%   weight is a free parameter, measured as not binding -- M-03), one derivative
%   chain per axis for a PDE, one chain per field for coupled fields, and the
%   constraint rows on chosen blocks (Neumann on the f'-block).

\subsection{Higher order, variable coefficients, and the padded block count}
\label{sec:app-variable-coeff}
% Moved here from the draft's Appendix B.

For an order-$k$ equation the state is $w = (\fhat{0};\, \dots;\, \fhat{k})$ and
$\Adesc$ stacks $k$ recurrence block-rows
$-\Dt\,\fhat{j} + \Bt\,\fhat{j+1} = 0$ above the single differential row
$\sum_j \M_{a_j}\,\fhat{j} = 0$: $k+1$ block-rows, each $N$ coefficients wide,
giving the classical $(k+1)N \times (k+1)N$ square matrix of
\Cref{def:descriptor-residual-operator}.

A variable coefficient $a_j(x)$ enters through $\M_{a_j}$, the square Galerkin truncation to the first $N$ modes of the exact product (the multiplication of \Cref{sec:pihm-residual}), built from its Chebyshev coefficients, or from its monomial coefficients converted exactly. $\M_{a_j}$ is banded, with bandwidth the degree of $a_j$, so it folds into the same shift-and-diagonal form as $\Bt$ and $\Dt$. The truncation is the only approximation in the linear construction; its error is confined to the top Chebyshev modes of a product that is smooth and resolved.

\paragraph{Padding to a power of two.} The classical $(k+1)N \times (k+1)N$
matrix above is exact and complete as a linear system; it is not, however,
the object any later stage of the pipeline actually builds. A gate-level
realisation must index which block $\fhat{j}$ a given operator acts on with a
qubit register, and a register of $m$ qubits addresses $2^m$ blocks, not
$k+1$ of them for an arbitrary $k$. $\Adesc$ is therefore \emph{embedded} in a
$2^{\lceil \log_2(k+1) \rceil}N$-dimensional space: the $k+1$ block-rows and
block-columns above are kept exactly as derived, and
$2^{\lceil \log_2(k+1) \rceil} - (k+1)$ further diagonal blocks are appended,
each a row $c_{\mathrm{pad}} I_N$ with $c_{\mathrm{pad}}=1$ pinning its own dummy field to zero,
$\fhat{j}_{\text{dummy}} = 0$; no other row reads a padding block, so the padding decouples. This padding is exact: the physical solution already satisfies the dummy rows
trivially, since it has nothing in those blocks, so no error is introduced
beyond the multiplication's truncation above. The qubit cost, however, is not
uniform in $k$:

\begin{center}
\begin{tabular}{ccccc}
\toprule
Order $k$ & Blocks $k+1$ & Field qubits $\lceil\log_2(k+1)\rceil$ & Padded blocks $2^{\lceil\log_2(k+1)\rceil}$ & Dummy blocks \\
\midrule
1 & 2 & 1 & 2 & 0 \\
2 & 3 & 2 & 4 & 1 \\
3 & 4 & 2 & 4 & 0 \\
4 & 5 & 3 & 8 & 3 \\
7 & 8 & 3 & 8 & 0 \\
\bottomrule
\end{tabular}
\end{center}

so the state $w$ the gate-level construction actually carries has
$2^{\lceil \log_2(k+1) \rceil}N$ components, not $(k+1)N$: equal to the
unpadded count exactly when $k+1$ is itself a power of two (orders $1, 3, 7,
\dots$), and up to nearly double it otherwise -- order $2$, the most common
case after order $1$, already pays for a field register it only three-quarters
uses. This is the qubit cost that must always accompany the gate-count
comparison of \Cref{sec:descriptor-cost}.

\subsection{Proof of \texorpdfstring{\Cref{thm:kernel-equivalence}}{kernel equivalence}}
\label{app:proof-kernel-equivalence}
% LAND: Bt is upper triangular with a non-zero diagonal, hence invertible; so
%   the recurrence rows hold iff f_hat^(j+1) = G f_hat^(j) for every j; then the
%   equation row of A_desc on (p, Gp, ..., G^k p) equals A_std p. Both
%   inclusions. The padding blocks vanish on the kernel.

\subsection{The term IR and its compilation}
\label{app:term-ir}
% LAND: the IR grammar (field word x shift x diagonal; diagonals as Z-strings
%   by eq:diag-index and reflections R_m); the banded-to-shift-diagonal
%   conversion (and its limit, implementation issue I-06); what alg:ir-to-lcu
%   compiles; how the IR mode of tab:verification-modes checks it against the
%   independently assembled operator at every n.

\subsection{Proof of \texorpdfstring{\Cref{prop:descriptor-cost}}{the descriptor cost}}
\label{app:proof-descriptor-cost}
% LAND: alpha_A = sum |c_t| (the dominant term zeta^j Dt, of norm 2 zeta^j (N-1), is
%   encoded at its norm, so alpha_A/||A_desc|| = 1 + O(1/N) to leading order on a regular
%   one-axis equation at fixed zeta; about 2 where the leading coefficient vanishes, since
%   comparable terms add in l1); the diagonal selector has ceil(log2(n+4)) qubits for R2
%   (n gates Z_j, the identity, three reflections), plus the field and shift selectors, the
%   overflow qubit and the projection flags: ceil(log2(n+4)) + 8 ancilla for R2's system rows
%   (n = 2-12); O(n^2) gates. Order-optimality among exact LCU encodings that expand
%   diagonals in Pauli strings from the Omega(log L) bound (chakraborty2025qsvt).

\subsection{Block rescaling}
\label{app:block-rescaling}
% LAND: the rescaling multiplies each LCU term's coefficient by zeta^m for its column block m:
%   no gate, no ancilla, only a different alpha; zeta is the a priori rule of
%   sec:descriptor-rescaling (never fitted); its effect on Delta and on alpha_H; why general
%   equilibration (Theta(dim w) gates, a diagonal block encoding) is excluded.
